\begin{conf-abstract}
{Noble gases as tracers for groundwater in a big high- altitude lake}
{Paul Moser Röggla, The ICDP Nam Core Team, Torsten Haberzettl, Hendrik Vogel, Rolf Kipfer}
{Eawag}
{Nam Co (30°42'N, 90°33'E; 4712 m a.s.l), consisting of a deeper main basin (max depth: 100m) and a shallower eastern basin (max depth: 50m), is one of the highest saline lakes in the world and is thought to exhibit substantial groundwater–lake interactions (see e.g. Ren et al., 2025). Recent mass balance estimates suggest that up to 25\% of the lake’s annual water input originates from lacustrine groundwater discharge (Ren et al., 2025). Terrestrial groundwater typically carries an excess air signature, characterized by oversaturation of atmospheric noble gases, while dissolved gases in surface water bodies are generally expected to remain in atmospheric equilibrium (Brennwald et al., 2013). However, noble gas measurements from the main basin of Nam Co reveal a distinct excess air signal at 50 m depth, with $\Delta$Ne oversaturations of up to 30\%. In addition, water at this depth shows a measurable tritium age, which is unexpected for a dimictic lake with seasonal overturn. We interpret these observations as direct measurable evidence for permafrost-derived groundwater discharge into the eastern basin, followed by lateral transport toward the central basin via horizontal, density-driven currents.

Brennwald, M.S. et al., In: Burnard, P. (Ed.), The Noble Gases as Geochemical Tracers. Springer, Berlin, Heidelberg, pp. 123–153. %http://dx.doi.org/10.1007/978-3-642-28836-4_6. 

Ren, Wenhao et al., Journal of Hydrology, 662, Part B, 2025, 133956%, https://doi.org/10.1016/j.jhydrol.2025.133956.
}
\end{conf-abstract}
